\par\Needspace{4\baselineskip}
\begin{center}
\SourceLeaf{024b}\pdfbookmark[2]{6.}{article-006}6.
\end{center}
We now turn to the relation between the motion and \emph{time}. As in Art. 1, we set the area swept out around the Sun in time $t$ $= \frac{1}{2}g$ and the mass of the moving body $= \mu$, taking the mass of the Sun $= 1$. We then have $g = kt\BareRoot{p}.\BareRoot{(1+\mu)}$. The differential of the area is $= \frac{1}{2}rr\,\mathrm{d}v$, which gives $kt\BareRoot{p}.\BareRoot{(1+\mu)}
= \int rr\,\mathrm{d}v$, with the integral taken so as to vanish for $t=0$. The integration must be carried out differently for the different kinds of conic sections. We shall therefore consider them separately, beginning with the ELLIPSE.

The fraction that determines $r$ from $v$ has two terms in its denominator. We shall first remove this inconvenience by introducing a new quantity in place of $v$. To this end we set $\operatorname{tang}\frac{1}{2}v\BareRoot{\frac{1-e}{1+e}}
= \operatorname{tang}\frac{1}{2}E$. The last formula of the preceding article then gives for $r$
\[
r = \frac{p\cos\frac{1}{2}E^2}{(1+e)\cos\frac{1}{2}v^2} = p\left(\frac{\cos\frac{1}{2}E^2}{1+e}+\frac{\sin\frac{1}{2}E^2}{1-e}\right) = \frac{p}{1-ee}(1-e\cos E).
\]
Furthermore, we obtain $\frac{\mathrm{d}E}{\cos\frac{1}{2}E^2} = \frac{\mathrm{d}v}{\cos\frac{1}{2}v^2}\BareRoot{\frac{1-e}{1+e}}$ and hence $\mathrm{d}v = \frac{p\,\mathrm{d}E}{r\BareRoot{(1-ee)}}$. It follows that
\[
rr\,\mathrm{d}v = \frac{r\,\grave{p}\,\mathrm{d}E}{\BareRoot{(1-ee)}} = \frac{pp}{(1-ee)^{\frac{3}{2}}}(1-e\cos E)\,\mathrm{d}E,
\]
and, on integrating,
\[
kt\BareRoot{p}.\BareRoot{(1+\mu)} = \frac{pp}{(1-ee)^{\frac{3}{2}}}(E-e\sin E)+\mathrm{Const.}
\]
If, then, we reckon time from the passage through perihelion, where $v=0$, $E=0$
\SourceLeaf{025}and hence the constant $= 0$, we obtain, since $\frac{p}{1-ee}=a$,
\[
E-e\sin E = \frac{kt\BareRoot{(1+\mu)}}{a^{\frac{3}{2}}}\cdot
\]

In this equation the auxiliary angle $E$, called the \emph{eccentric anomaly}, must be expressed in radians. We may of course continue to express this angle in degrees and their subdivisions, provided that $e\sin E$ and $\frac{kt\BareRoot{(1+\mu)}}{a^{\frac{3}{2}}}$ are expressed in the same way. These quantities will be expressed in seconds of arc if multiplied by $206264{,}806$. We can dispense with multiplying the latter quantity by using $k$ already expressed in seconds of arc and thus substituting $k = 3548^{\prime\prime}{,}18761$ for the value given above. Its logarithm is $= 3{,}550\,0065\,746$. — Expressed in this way, the quantity $\frac{kt\BareRoot{(1+\mu)}}{a^{\frac{3}{2}}}$ is called the \emph{mean anomaly}. It therefore increases in proportion to time, with a daily increment $\frac{k\BareRoot{(1+\mu)}}{a^{\frac{3}{2}}}$, which is called the \emph{mean daily motion}. We shall denote the mean anomaly by $M$.
